% From mitthesis package
% Version: 1.01, 2023/06/19
% Documentation: https://ctan.org/pkg/mitthesis
%
% The abstract environment creates all the required headers and footnote. 
% You only need to add the text of the abstract itself.
%
% Approximately 500 words or less; try not to use formulas or special characters
% If you don't want an initial indentation, do \noindent at the start of the abstract

The continuous demand for higher performance, energy efficiency, and scalability in information processing 
systems has driven integrated technologies beyond the limits of purely electronic solutions. Although 
electronic integrated circuits have enabled major advances in computing and communication, fundamental 
constraints related to bandwidth, power consumption, and interconnect scalability have become increasingly 
significant as system complexity grows. In this context, integrated photonics has emerged as a key enabling 
technology, exploiting the properties of light to overcome these limitations. Early photonic integrated 
circuits were predominantly developed following an application-specific design paradigm, which enabled 
highly efficient and compact solutions but resulted in long development cycles, limited flexibility, and 
poor adaptability to evolving system requirements. Motivated by the success of programmable electronic 
hardware, programmable photonic integrated circuits have been proposed as a promising alternative, offering 
software-controlled reconfiguration on a common hardware platform, yet their widespread adoption remains 
constrained by challenges related to scalability, reliability, and system-level integration.

The primary objective of this doctoral thesis is to address the scalability challenges of programmable 
photonic integrated circuits across multiple dimensions. At the device and circuit levels, the thesis 
investigates design strategies that enable the integration of a large number of tunable elements on a 
single chip while maintaining reliable and stable operation. Particular emphasis is placed on improving 
fabrication tolerance and mitigating the impact of process variations, environmental fluctuations, 
and long-term drift, which become increasingly critical as circuit complexity increases. By enhancing 
robustness at the design stage, the proposed approaches aim to reduce calibration overhead and improve 
reproducibility in large-scale programmable photonic systems.

At the architectural and system levels, this work explores scalable programmable photonic architectures 
with efficient use of on-chip area and reduced control complexity. The thesis also addresses practical 
challenges associated with electrical and optical packaging, including power consumption, thermal management, 
and interconnection density, which often represent key bottlenecks in large-scale programmable PICs. 
Through the co-design of photonic architectures and their associated control and packaging strategies, 
this research seeks to enable realistic and deployable programmable photonic platforms.

Finally, the thesis demonstrates the applicability of the proposed scalable programmable PICs through 
experimental validation and representative application scenarios. These demonstrations highlight the 
potential of scalable, reliable, and fabrication-tolerant programmable photonic circuits in areas such as 
optical signal processing, communications, and emerging photonic computing paradigms. Overall, this work 
contributes to the advancement of programmable photonic integrated circuits from flexible research 
demonstrators toward robust and scalable technologies capable of meeting the demands of future photonic 
systems.

% The field of programmable silicon photonics has garnered significant attention in recent years due to its potential to revolutionize various technological domains.
% The ability to reconfigure photonic integrated circuits on demand offers unprecedented flexibility and adaptability for applications ranging from telecommunications and optical signal processing to sensing and quantum information processing.
% This interest is driven by the promise of creating versatile photonic platforms that can be tailored to specific tasks through software control, much like their electronic counterparts.

% This manuscript provides a comprehensive exploration of the theoretical underpinnings, the developmental journey, and the experimental validation of applications stemming from this research.
% The central theme revolves around the design, implementation, and application of a multipurpose programmable photonic integrated circuit (PIC) based on a recirculating mesh architecture.
% The thesis outlines the critical need for a robust technology stack, with a particular emphasis on the software infrastructure required to harness the full potential of photonic integrated circuits.
% A significant portion of this work is dedicated to the in-depth documentation of various applications demonstrated on the developed photonic processor, including advancements in optical networking, tunable delay lines, optical signal processing, and topological photonics, enabled by specifically developed algorithms and high-level instructions.
% Furthermore, the thesis delves into the theoretical intricacies of optical arbitrary matrix-multiply operators, proposing a novel theoretical model for Mach-Zehnder interferometers within a recirculating mesh array.
% This theoretical framework is complemented by the software implementation and experimental demonstration of arbitrary parametric matrix-multiply operators on the fabricated photonic processor.
% Finally, the thesis concludes by discussing the key findings, the challenges encountered and addressed, and the commercial and research prospects of the developed technology, while also outlining future directions to further the impact and accessibility of multipurpose programmable photonic platforms.
